\documentclass[a4paper,12pt]{article}
\usepackage[utf8]{inputenc}
\usepackage[russian]{babel}
\usepackage{amsmath, amssymb}

\begin{document}

\section{Построение решётки по проверочной матрице}

Рассмотрим заданную проверочную матрицу линейного кода:

\[
G =
\begin{pmatrix}
1 & 0 & 1 & 1 & 0 & 1 \\
1 & 0 & 1 & 0 & 1 & 0 \\
1 & 1 & 0 & 1 & 0 & 0
\end{pmatrix}
\]

Код является двоичным линейным с параметрами $[6, 3]$.

Синдромная решётка строится следующим образом:

\begin{enumerate}
    \item Ярусы $0, 1, \dots, 6$ соответствуют этапам обработки $i$ символов, где $i = 0, \dots, 6$.
    
    \item Состояния на $i$-м ярусе представляют собой все возможные частичные синдромы
    \[
    s_i = G_{[:, 0:i]} \cdot r_{0:i}^T \pmod{2}.
    \]
    
    \item Из каждого состояния $s$ на ярусе $i-1$ исходят две ветви:
    \begin{itemize}
        \item при метке $0$: новое состояние $s' = s$;
        \item при метке $1$: новое состояние $s' = s + g_i$,
    \end{itemize}
    где $g_i$ — соответствующий столбец матрицы.
\end{enumerate}

Достижимые состояния (в формате $(s_1, s_2, s_3)$) на каждом ярусе:

\begin{itemize}
    \item Ярус 0: $\{(0,0,0)\}$ (одно состояние)
    
    \item Ярус 1: $\{(0,0,0), (1,1,1)\}$ (два состояния)
    
    \item Ярус 2: $\{(0,0,0), (0,0,1), (1,1,0), (1,1,1)\}$ (четыре состояния)
    
    \item Ярус 3: $\{(0,0,0), (0,0,1), (1,1,0), (1,1,1)\}$ (четыре состояния)
    
    \item Ярус 4: все восемь возможных синдромов:
    \[
    \{(0,0,0), (0,0,1), (0,1,0), (0,1,1), (1,0,0), (1,0,1), (1,1,0), (1,1,1)\}
    \]
    
    \item Ярус 5: также все восемь возможных синдромов
    
    \item Ярус 6: все восемь возможных синдромов
\end{itemize}

На последнем ярусе ($i = 6$) корректные кодовые слова завершаются исключительно в состоянии $(0,0,0)$.

\section{Совпадение с синдромной решёткой}

Построенная структура по определению является синдромной решёткой данного кода, поскольку её состояния соответствуют частичным синдромам
\[
s_i = \sum_{j=1}^{i} r_j \cdot h_j,
\]
а переходы между состояниями определяются столбцами матрицы $G$.

Следовательно, полученная решётка полностью совпадает с синдромной решёткой. Это подтверждается непосредственным построением: на каждом уровне множество достижимых состояний совпадает с множеством всех возможных синдромов, возникающих после обработки первых $i$ символов.

\end{document}